\documentclass[11pt,letterpaper]{article}
\usepackage[letterpaper,margin=1in,headheight=14pt]{geometry}
\usepackage{mathptmx}
\usepackage[T1]{fontenc}
\usepackage[utf8]{inputenc}
\usepackage{microtype}
\usepackage{amssymb}
\usepackage{booktabs}
\usepackage{tabularx}
\usepackage{longtable}
\usepackage{array}
\usepackage{xcolor}
\usepackage{enumitem}
\usepackage{fancyhdr}
\usepackage[most]{tcolorbox}
\usepackage{titlesec}
\usepackage[hidelinks]{hyperref}

\definecolor{deep}{HTML}{1F3A5F}
\definecolor{warn}{HTML}{9C3B2E}
\definecolor{soft}{HTML}{F2F4F7}

\titleformat{\section}{\normalfont\Large\bfseries\color{deep}}{\thesection.}{0.5em}{}
\titleformat{\subsection}{\normalfont\large\bfseries}{\thesubsection.}{0.5em}{}
\titlespacing*{\section}{0pt}{2.3ex}{1.1ex}
\pagestyle{fancy}\fancyhf{}
\lhead{\small\color{deep}Solutions Dimension --- Survey Instrument}
\rhead{\small\thepage}
\renewcommand{\headrulewidth}{0.4pt}
\setlist{itemsep=2pt,parsep=0pt,topsep=3pt}
\renewcommand{\arraystretch}{1.22}
\newcolumntype{R}{>{\raggedright\arraybackslash}X}

\newtcolorbox{note}[1]{colback=soft,colframe=deep!45,fonttitle=\bfseries,
  title={#1},boxrule=0.6pt,arc=2pt,left=6pt,right=6pt,top=4pt,bottom=4pt}
\newtcolorbox{caution}[1]{colback=warn!6,colframe=warn,fonttitle=\bfseries,
  title={#1},boxrule=0.8pt,arc=2pt,left=6pt,right=6pt,top=4pt,bottom=4pt}

\newcommand{\R}{\textsuperscript{R}}
\newcommand{\anchor}{\textsuperscript{\dag}}

\begin{document}

\begin{titlepage}
\centering
\vspace*{2.4cm}
{\color{deep}\rule{\textwidth}{2pt}}\\[0.8em]
{\Huge\bfseries\color{deep} Survey Instrument}\\[0.6em]
{\LARGE The Solutions Dimension}\\[0.4em]
{\large Participant and facilitator forms, with scoring notes}\\[0.5em]
{\color{deep}\rule{\textwidth}{2pt}}\\[2.4cm]
\begin{minipage}{0.84\textwidth}
\begin{note}{What this instrument is for}
It measures three things at once: whether participants found the activity usable
and worthwhile, whether the four design requirements from Cycle 1 were actually
met, and whether confidence on specific tasks moved. Items marked \anchor{} are
carried verbatim from the Cycle 1 survey so the two rounds stay comparable.
Items marked \R{} are reverse-scored.
\end{note}
\end{minipage}
\vfill
{\small Version 1.0 \quad$\cdot$\quad Draft for IRB review}
\end{titlepage}

\tableofcontents
\newpage

% =====================================================================
\section{Before You Administer This}

\begin{caution}{The one mistake to avoid}
Satisfaction is not learning. Cycle 1 reported that participants enjoyed the
activity and found it easy, and reviewers correctly objected that this said
nothing about whether anyone learned anything. This instrument has the same
limit. It can support claims about \emph{usability, perceived value, workload,
and confidence}. It cannot support a claim about learning gains --- only the
performance measures M1--M4 in the study design can do that, and they come from
the session artifacts, not from here.

Write the claims down before you collect anything, and check afterwards that
none of them drifted.
\end{caution}

\subsection{Length and Timing}

The whole participant instrument is 38 items and runs 6--8 minutes. That is
already longer than a single play session, so resist adding to it. If you must
cut, drop Part F (cognitive load) first and Part D (realism) second; keep Parts
B, C and E intact.

\begin{center}
\small
\begin{tabularx}{\textwidth}{@{}l l R@{}}
\toprule
\textbf{When} & \textbf{What} & \textbf{Why then} \\
\midrule
Before play, 2 min & Part A (background), Part B (pre-confidence) & Confidence
must be measured before exposure or the pre/post comparison is worthless \\
After the debrief, 6 min & Parts C--H & After debrief, not after Phase 3 --- the
debrief is part of the intervention \\
Same day & Facilitator form (Part I) & Memory for interventions decays fast \\
\bottomrule
\end{tabularx}
\end{center}

\subsection{Linking Pre to Post Without Names}

Participants generate their own code, which lets you pair pre and post responses
without collecting identity. Put this at the top of both forms:

\begin{quote}\itshape
Please create your code: the first two letters of the street you grew up on, the
last two digits of your phone number, and the first two letters of your mother's
first name. For example, ELM Street, phone ending 47, mother Rosa gives
\texttt{EL47RO}. Use exactly the same code on both surveys.
\end{quote}

Expect 10--15\% of codes not to match. Report the pairing rate.

\subsection{Response Format}

All agreement items use a five-point scale: \textbf{1 Strongly disagree,
2 Disagree, 3 Neither agree nor disagree, 4 Agree, 5 Strongly agree.} Confidence
items use a 0--100 slider. Cognitive load items use 0--10.

Do not offer a ``not applicable'' option on agreement items; it invites
satisficing. Do allow every item to be skipped.

% =====================================================================
\section{Part A --- Background}

Reviewers of the Cycle 1 paper specifically asked for this. Six items.

\begin{enumerate}[label=A\arabic*.]
\item What is your primary field of study? \emph{(Computer science / Information
technology / Engineering / Natural sciences / Mathematics / Social sciences /
Humanities / Business / Health sciences / Other)}
\item What is your current level? \emph{(First or second year undergraduate /
Third or fourth year undergraduate / Graduate student / Faculty or staff /
Other)}
\item How many cybersecurity courses have you completed? \emph{(None / One /
Two or three / More than three)}
\item Before today, had you heard of the CIS Critical Security Controls?
\emph{(No / Heard of them / Used them)}
\item How often do you play board or card games with other people?
\emph{(Never / A few times a year / Monthly / Weekly or more)}
\item In your team today, which best describes you? \emph{(I did most of the
talking / I contributed about equally / I mostly listened / Unsure)}
\end{enumerate}

\begin{note}{Why A5 and A6 are here}
A5 tests a confound worth ruling out: if only experienced gamers enjoyed it, the
activity is less useful than it looks. A6 gives you a cheap check on whether
non-majors were talked over, which is the failure mode the facilitator manual
warns about, and it is more honest than asking the facilitator to judge it.
\end{note}

% =====================================================================
\section{Part B --- Confidence, Before Play}

Administer \emph{before} the session. Same five items reappear as Part G.

\begin{quote}\small
How confident are you that you could do each of the following right now?
Slide from 0 (cannot do at all) to 100 (highly certain I can do).
\end{quote}

\begin{enumerate}[label=B\arabic*.]
\item Explain to a non-technical person why an organization should keep backups.
\item Given a described attack, name two security measures that would make it
less likely to succeed.
\item Choose between two security measures when the organization can only afford
one, and say why.
\item Explain the difference between making an attack less likely and making it
less damaging.
\item Tell an organization what risk remains after the measures they have put in
place.
\end{enumerate}

\begin{note}{Why these five}
They are written to Bandura's guidance for self-efficacy scales: specific tasks
rather than general traits, and graded confidence rather than agreement. Each
maps to something the game actually asks a player to do --- B3 is the budget,
B4 is the risk matrix, B5 is Phase 3 --- so movement is interpretable rather
than diffuse.
\end{note}

% =====================================================================
\section{Part C --- The Cycle 1 Anchor Items}

Carry these four verbatim. They are the only bridge between the two rounds, and
rewording them for elegance would break it.

\begin{enumerate}[label=C\arabic*.]
\item\anchor{} The Security Cards helped me learn.
\item\anchor{} The Security Cards are easy to use.
\item\anchor{} The Security Cards raised or maintained my cybersecurity awareness.
\item\anchor{} The Security Cards were not a source of frustration.
\end{enumerate}

\begin{note}{Reporting these}
Report counts, not percentages. With nineteen or thirty participants,
``31.6\%'' claims a precision the sample does not have; ``6 of 19'' is honest and
easier to read. A diverging stacked bar beats four pie charts.
\end{note}

% =====================================================================
\section{Part D --- Content of the Cards}

Tests design goal G1 (vocabulary accessibility), G3 (framework traceability),
and requirement R2 (bound the working set).

\begin{enumerate}[label=D\arabic*.]
\item The cards used language I could understand without asking for help.
\item The examples on the back of the cards helped me decide what a card meant.
\item\R{} There were too many cards to keep track of.
\item\R{} Some cards seemed to overlap with each other.
\item Seeing that a card referred to a published control, such as ``CIS Control
7,'' made the card feel more trustworthy.
\item If I wanted to learn more about one of these controls after today, I would
know where to start.
\end{enumerate}

\begin{note}{D3 and D4 are the honest ones}
Both are direct quotations of Cycle 1 complaints, reworded as items. If the
tiering and the eight-card hand worked, D3 should fall relative to Cycle 1's
open-ended reports. If it does not, say so --- a requirement that was not met is
a finding, not a failure.
\end{note}

% =====================================================================
\section{Part E --- Process and Gameplay}

Tests R1 (choices must be contestable) and R3 (basis for feedback on quality).
E3 is the single most important item in this instrument.

\begin{enumerate}[label=E\arabic*.]
\item The order of the activities was clear to me.
\item Having a limited budget made the choices harder in a way that was useful.
\item\R{} There was usually one obviously correct answer.
\item My team disagreed about at least one purchase.
\item Having to justify a modifier card out loud made me think more carefully.
\item I had a sense of whether our choices were any good.
\item We had enough time to finish without rushing.
\item\R{} The facilitator had to explain the rules more often than I expected.
\end{enumerate}

\begin{note}{Reading E3 and E4 together}
Cycle 1's core complaint was that too many answers seemed equally acceptable.
E3 is that complaint stated as an item and E4 is its behavioural counterpart. If
the budget mechanism did what it was designed to do, E3 falls and E4 rises. If
E3 stays high, the budget is too loose --- which is a calibration problem with a
known fix, not a refutation of the design.

E6 will probably score low. Novices want an answer key and the activity
deliberately withholds one. Report it and interpret it rather than treating it
as a defect.
\end{note}

% =====================================================================
\section{Part F --- Realism and Transfer}

\begin{caution}{A limit worth stating in the paper}
Novices are poor judges of whether a security scenario is realistic --- they have
no basis for comparison, which is why they are in the course. Treat Part F as
measuring \emph{perceived relevance and intent to transfer}, not realism.
Realism is judged on the facilitator form by someone with domain knowledge, and
that is the number to cite if you make a claim about it.
\end{caution}

\begin{enumerate}[label=F\arabic*.]
\item The organization in the scenario felt like a place that could exist.
\item The trade-offs felt like ones a real organization would actually face.
\item I can see how this connects to advising an actual client.
\item Placing the threat again after buying controls matched how I think risk
really behaves.
\item\R{} This felt more like a puzzle than like a real decision.
\item I would use something I learned today if a friend or family member asked me
a security question.
\end{enumerate}

% =====================================================================
\section{Part G --- Confidence, After Play}

Repeat items B1--B5 verbatim, same 0--100 slider, same wording. Label them
G1--G5. Do not show participants their earlier answers.

% =====================================================================
\section{Part H --- Workload, Features, and Open Responses}

\subsection{Workload (0--10)}

Three items, separating load that comes from the subject from load that comes
from how it was presented. The distinction matters: high intrinsic load is
expected and fine, high extraneous load is a design defect you can fix.

\begin{enumerate}[label=H\arabic*.]
\item The topic itself --- threats and security controls --- was complex.
\emph{(0 = very easy, 10 = very complex)}
\item The way the activity was set up and explained made it harder than it
needed to be. \emph{(0 = not at all, 10 = very much)}
\item The activity helped me understand how threats and controls relate.
\emph{(0 = not at all, 10 = very much)}
\end{enumerate}

\subsection{Which Parts Were Worth It}

Best-worst rather than rating everything, because rating eight features on a
five-point scale produces eight fours and tells you nothing.

\begin{quote}\small
Of the parts of the activity listed below, which \textbf{one} was most useful to
you, and which \textbf{one} was least useful?
\end{quote}

\begin{enumerate}[label=(\alph*)]
\item The four threat cards used to build the scenario
\item The Solutions cards describing security measures
\item The CIS control references printed on the cards
\item The risk matrix
\item The limited budget and spending tokens
\item The modifier cards and having to justify them
\item Placing the threat a second time after buying controls
\item Comparing your team's choices with another team's
\end{enumerate}

\begin{note}{Scoring best-worst}
For each option, score $= (\text{times chosen most useful}) - (\text{times chosen
least useful})$, divided by the number of respondents. The result runs from $-1$
to $+1$ and ranks cleanly. Report the whole ranking, including the negatives ---
a feature that most participants found least useful is a candidate to cut, and
saying so is more credible than reporting only what people liked.
\end{note}

\subsection{Open Responses}

Three, all specific. Generic prompts such as ``any other comments'' produce
generic answers.

\begin{enumerate}[label=H\arabic*.,start=4]
\item Which purchase did your team argue about most, and how did you settle it?
\item What did you want to buy and could not afford? What would you have told the
organization about that gap?
\item If you could change one thing about this activity, what would it be?
\end{enumerate}

% =====================================================================
\section{Part I --- Facilitator Form}

One per facilitator per session, completed the same day. This is where domain
judgment lives and where the process data that no participant can report gets
captured.

\begin{enumerate}[label=I\arabic*.]
\item Scenario used, budget issued, gameplay variant, number of teams, team sizes.
\item Actual elapsed minutes per phase, and total including debrief.
\item As someone with security background: how realistic was this scenario for
the kind of organization it described? \emph{(1--5, plus one line of comment)}
\item Did teams argue substantively about at least one purchase?
\emph{(All teams / Most / Some / None)}
\item Roughly how many interventions did you make per table, and at which rung?
\emph{(Counts for rungs 1--5 from the facilitator manual)}
\item Did any of the four immediate-intervention situations arise? \emph{(Yes/No
per situation; do not record details of disclosures or personal content)}
\item Justification quality, scored 0--3 per team using the manual's rubric.
\item What broke, confused people, or took longer than expected?
\item What would you change before running it again?
\end{enumerate}

\begin{caution}{On I6}
Record only that a situation occurred, never its content. If a participant
disclosed a real vulnerability or a scenario touched personal experience, that
belongs in your own follow-up, not in a research dataset. Say this explicitly in
your IRB protocol.
\end{caution}

% =====================================================================
\section{Scoring and Analysis Notes}

\subsection{Reverse-Scored Items}

D3, D4, E3, E8, F5. Recode as $6 - x$ before combining with anything.

\subsection{Suggested Groupings}

These are provisional. With a first sample you can report item-level results and
compute internal consistency; do not present them as validated subscales until
you have the data to justify it.

\begin{center}
\small
\begin{tabularx}{\textwidth}{@{}l l R@{}}
\toprule
\textbf{Grouping} & \textbf{Items} & \textbf{What it speaks to} \\
\midrule
Accessibility & D1, D2, D6, C2 & Goal G1, and the Cycle 1 ease-of-use result \\
Cognitive burden & D3\R, D4\R, H2 & Requirement R2 \\
Contestability & E2, E3\R, E4 & Requirement R1 --- the central design claim \\
Feedback and rigour & E5, E6 & Requirement R3 \\
Traceability & D5, D6 & Goal G3 \\
Perceived transfer & F2, F3, F6 & Whether it connects to clinic work \\
Confidence change & G1--G5 minus B1--B5 & Paired, per participant \\
\bottomrule
\end{tabularx}
\end{center}

\subsection{Analysis}

Report counts and medians for ordinal items rather than means alone; with samples
this size a mean of 4.42 and a mean of 4.11 are not meaningfully different.
Use Wilcoxon signed-rank for the paired confidence change, report an effect size
with a confidence interval, and treat the study as powered only for large
effects. State the pre/post pairing rate.

\subsection{What You Can and Cannot Say}

\begin{center}
\small
\begin{tabularx}{\textwidth}{@{}R R@{}}
\toprule
\textbf{Supported by this instrument} & \textbf{Not supported} \\
\midrule
Participants found the deck usable and its vocabulary accessible &
Participants learned to select appropriate controls \\
The budget made choices feel contestable &
The budget improved decision quality \\
Confidence on five specific tasks rose after one session &
Competence rose, or the change persisted \\
Participants rated the trade-offs as plausible &
The scenarios are realistic \emph{(use the facilitator judgment, I3)} \\
Which features participants valued most &
Which features contributed most to learning \\
\bottomrule
\end{tabularx}
\end{center}

The right-hand column is answerable, but from the session artifacts --- the
exported portfolios, matrix deltas, and justification scores --- not from a
questionnaire.

\end{document}
